% !TEX root = ./main.tex
% Estrutura exigida pela PRPPG: Introdução, Materiais e Métodos,
% Resultados e Discussão, Conclusões/Considerações finais, Referências.
% Entre 4 e 6 mil palavras ou 10 a 15 páginas (sem as referências).

\section*{Introdução}

Problemática, objetivos e justificativa da pesquisa \parencite{doCarmo1992}.

\section*{Materiais e Métodos}

Materiais (fontes, bibliografia base) e metodologia da pesquisa
\parencite{lima1973variedades}.

\section*{Resultados e Discussão}

\subsection*{Título da Subseção}

\begin{definition}
  \label{def:exemplo}
  Definição de exemplo.
\end{definition}

\begin{theorem}
  \label{teo:exemplo}
  Pela \cref{def:exemplo}, vale \ldots
\end{theorem}

\begin{proof}
  Demonstração.
\end{proof}

\section*{Conclusões}

Considerações finais (sem citações e sem assuntos novos).

\printbibliography[title={Referências}]
